\begin{abstract}
A decode server's batch is capped by device memory, not arithmetic: KV cache
exhausts HBM long before the arithmetic units are busy, so the GPU runs at
single-digit utilization while turning requests away. Coherent CPU--GPU
interconnects make host DRAM fast enough to hold the KV of an \emph{actively
decoding} request, but existing systems vary how much KV sits on the host and
never \emph{whose} --- they place none, or a uniform fraction of every request.

We present a heterogeneous KV memory that manages HBM and host DRAM as one
store and assigns each request to a tier. A \emph{host-tier} request holds no
device blocks at all: its KV lives in host memory and is streamed in one layer
at a time, just before the attention that consumes it, through a small buffer
that is reused across layers --- so it occupies a few percent of the device
memory a resident request needs. Both tiers are served in a single forward
pass, so the weights are read once however the batch is split, and an admission
controller sizes the host cohort from the link bandwidth and step time it
measures at runtime.

A capacity model then says which configurations this helps. It takes a
profile of the device --- the bandwidth each read achieves, the step time, and
the two costs the host path adds --- and returns the host cohort that maximises
decode throughput. Its content is that the value of a second tier tracks
\textbf{how much of memory the weights consumed}, not model size and not
context length, which cancels.

We implement this in vLLM and measure the costs it turns on: a host-tier
request adds a fixed $12.3$\,ms per step plus $2.2$\,ms per request, so a
prefetch one layer deep hides about $90\%$ of each transfer and a per-step
constant decides the rest. Applied to a Grace-Hopper node the model projects
$\times1.50$ and $\times1.67$ more served requests for the two
weight-heavy configurations, and throughput gains of $\times1.01$ and
$\times1.05$ against margins of a few milliseconds in that constant --- which
is the measurement we ask of such a node, and the one PCIe cannot supply.
\end{abstract}
